\documentclass[12pt]{article}
\usepackage{amsmath}

\usepackage{xcolor}

\usepackage{listings}
\usepackage[most]{tcolorbox}
\usepackage[colorlinks, allcolors=blue]{hyperref}


\title{E.B.1.1 - FOL: studenti ansiosi}
\author{}
\date{}

\begin{document}
\maketitle
\vspace{-5em}

\begin{quote}
``Nessuno studente supererà un qualunque esame se è ansioso o non ne studia il programma.''
\end{quote}

\begin{itemize}
    \item si può interpretare un predicato $Studia(x, y)$ come vero se lo studente $x$ studia il programma dell'esame $y$ e scrivere una modellazione di questo tipo:
\begin{align*}
    \forall x \forall y (& \bigr(Studente(x) \land Esame(y) \land \\
                       &(Ansioso(x) \lor \neg Studia(x, y))\bigr) \to \neg Supera(x, y) )
\end{align*}

    \item se si vuole invece modellare il programma come un'entità separata dall'esame, si può scrivere:
\begin{align*}
    \forall x \forall y \forall z (& \bigr(Studente(x) \land Esame(y) \land Programma(y, z)\\
                                   &\land (Ansioso(x) \lor \neg Studia(x, z))\bigr) \to \neg Supera(x, y) )
\end{align*}
\item (nel caso sopra un esame può anche avere più di un programma (e.g. sezioni del corso), e basta non studiarne uno per non passare - se invece assumiamo che ogni esame abbia esattamente un programma possiamo modellare ``programma'' come simbolo di funzione)
\end{itemize}

Abbiamo quindi:
$$\mathcal{P} = \{ Studente/1, Esame/1, Programma/2, Ansioso/1, Studia/2, Supera/2 \}$$

(in cui i significati dei predicati sono quelli intuitivi)

\begin{itemize}
    \item un primo modello (t.c. $M_1 \models \phi$) è $M_1 = \langle D, I \rangle$, dove:
\begin{itemize}
    \item $D = \{ st, es, pr \}$
    \item $I(Studente) = \{ st \}$
    \item $I(Esame) = \{ es \}$
    \item $I(Programma) = \{ (es, pr) \}$
    \item $I(Ansioso) = \{ st \}$
    \item $I(Studia) = \emptyset$ 
    \item $I(Supera) = \emptyset$
\end{itemize}

\item un secondo modello (t.c. $M_1 \not\models \phi$) è $M_2 = \langle D, I \rangle$, dove:  

\begin{itemize}
    \item $D = \{ st, es, pr \}$
\item $I(Studente) = \{ st \}$
\item $I(Esame) = \{ es \}$
\item $I(Programma) = \{ (es, pr) \}$
\item $I(Ansioso) = \{ st \}$
\item $I(Studia) = \{ (st, pr)  \}$
\item $I(Supera) = \{ (st, es) \}$
\end{itemize}

\end{itemize}
\end{document}
